% !TeX root = ../main.tex
\section{Lexical Analysis}

This is where the pipeline actually starts. Everything in the previous block was about
languages as sets of strings, recognized by grammars and automata. Now we build the first
real piece of a language implementation: the component that reads raw text and turns it
into something a parser can work with.

\subsection{what a lexer does}

A \textbf{lexer} (also called a \emph{scanner} or \emph{tokenizer}) takes a raw stream of
characters and groups it into a stream of \textbf{tokens}, discarding anything that isn't
meaningful --- whitespace, comments, the exact number of spaces you used to indent a line.
It's the first phase of any compiler or interpreter, and everything downstream (parsing,
semantic analysis, code generation) works on tokens, never on raw characters again.

\begin{definition}[Token]
A \emph{token} is a pair $(\mathit{kind}, \mathit{lexeme})$, where $\mathit{kind}$ is the
category of the token (e.g.\ \texttt{NUMBER}, \texttt{IDENT}, \texttt{PLUS}) and
$\mathit{lexeme}$ is the exact substring of the source text that produced it. Two tokens
of the same kind can have different lexemes: \texttt{NUMBER(3)} and \texttt{NUMBER(42)}
are both \texttt{NUMBER} tokens.
\end{definition}

Concretely, given the input
\begin{lstlisting}
x = 3 + 42
\end{lstlisting}
a lexer produces something like
\[
\texttt{IDENT(x)} \;\; \texttt{ASSIGN(=)} \;\; \texttt{NUMBER(3)} \;\;
\texttt{PLUS(+)} \;\; \texttt{NUMBER(42)} \;\; \texttt{EOF}
\]
Five tokens from eleven characters. The spaces are gone, and so is any trace that they
were ever there --- which is exactly the point. The parser that consumes this stream
doesn't care how many spaces separated \texttt{x} and \texttt{=}; it only cares about the
sequence of token kinds.

\subsection{tokens are regular languages, lexers are DFAs}

This is the connection back to the previous block, and it's not a metaphor: it's literally
the same machinery. Each token category is defined by a regular expression --- identifiers
match \texttt{[a-zA-Z\_][a-zA-Z0-9\_]*}, integer literals match \texttt{[0-9]+}, the keyword
\texttt{if} matches the literal string \texttt{if}. A regular expression can be compiled
into a DFA (via Thompson's construction and then subset construction, if you did that
exercise in the previous block), and a DFA can decide in one linear pass over the input
whether a string belongs to the language it recognizes.

A lexer is what you get when you run \emph{all} of these DFAs at once, character by
character, and at each point ask ``which of these automata are still alive, and which one
(if any) is in an accepting state right now?'' Lexer generators like \texttt{lex}/\texttt{flex}
build exactly this: they take a list of regular expressions, one per token kind, and
compile them into a single combined DFA via a product construction. Wilhelm and Seidel's
\emph{Compiler Design: Syntactic and Semantic Analysis} works through this construction in
detail if you want the automaton-theoretic view all the way down; Aho, Lam, Sethi and
Ullman's \emph{Compilers: Principles, Techniques, and Tools} (the ``Dragon Book'') covers
the same ground in its lexical analysis chapter.

\subsection{maximal munch}

There's one rule that makes all of this work when several DFAs could match a prefix of
the input: \textbf{longest match}, also called \textbf{maximal munch}.

\begin{definition}[Maximal munch]
When more than one prefix of the remaining input matches some token pattern, the lexer
picks the \emph{longest} one. If two patterns match the same longest prefix, the pattern
listed first (usually: keywords before identifiers, specific operators before their
prefixes) wins.
\end{definition}

The classic example is why \texttt{<=} is one token and not two. Suppose your language has
both \texttt{<} (less-than) and \texttt{<=} (less-or-equal). When the lexer sees the
character \texttt{<}, it does \emph{not} immediately emit an \texttt{LT} token. It looks
one character ahead: if the next character is \texttt{=}, both characters together form a
longer match (\texttt{<=}), so that's what gets emitted. Get this wrong and
\texttt{x <= 10} tokenizes as \texttt{LT}, \texttt{ASSIGN}, \texttt{NUMBER} --- which is not
even close to what the program means.

The same principle governs identifiers versus keywords: \texttt{iffy} is one identifier
token, not the keyword \texttt{if} followed by the identifier \texttt{fy}. Maximal munch
says: keep consuming characters as long as the identifier pattern keeps matching. Only once
you \emph{cannot} extend the match any further do you ask, ``is this exact lexeme one of the
reserved keywords?''

\subsection{keywords are identifiers with a special case}

This gives the standard recipe for handling keywords, and it's worth stating explicitly
because it's easy to get backwards. You do \emph{not} write a separate regular expression
per keyword and try to match those first. Instead:

\begin{enumerate}
  \item Match the longest possible identifier-shaped lexeme (letters, digits, underscores,
    not starting with a digit).
  \item Look the resulting lexeme up in a fixed table of reserved words.
  \item If it's in the table, emit the corresponding keyword token
    (\texttt{IF}, \texttt{WHILE}, \texttt{PRINT}, \ldots). Otherwise, emit \texttt{IDENT}.
\end{enumerate}

This is simpler than it sounds, and it composes correctly with maximal munch: since
identifiers and keywords share exactly the same character pattern, there is never a
question of which pattern is ``longer.'' You match once, then classify.

\subsection{whitespace, comments, and position tracking}

Whitespace and comments are recognized and then thrown away; they never become tokens.
A lexer typically has a small loop that runs before scanning each token: skip spaces,
tabs, and newlines; if you see a comment marker, skip to the end of the comment (end of
line, for a line comment; the matching closer, for a block comment) and then go back to
skipping whitespace, since comments and whitespace can alternate.

Newlines deserve special mention: even though they're discarded like other whitespace,
the lexer must \emph{count} them, along with columns, to support error messages like
\texttt{line 4, column 12: unexpected character '\$'}. Without position tracking, every
error message you produce is useless for anyone actually trying to fix their program ---
imagine a 200-line file and the only feedback is ``syntax error somewhere.''

\subsection{lexical errors}

A lexical error happens when the current character doesn't start \emph{any} valid token
pattern --- an illegal character. In most languages this is things like a stray \texttt{\$}
or \texttt{\textbackslash} that isn't part of any operator, or an unterminated string
literal that runs off the end of the file before its closing quote. Note the scope of what
a lexer can catch: it only knows about individual characters and local patterns. It has no
idea whether \texttt{x} was declared, or whether \texttt{3 + true} type-checks. Those are
problems for later phases (parsing catches structural nonsense; semantic analysis, covered
later in this block, catches everything context-dependent).

\subsection{a hand-written lexer}

Let's build one. The toy language has integers, identifiers, the keywords \texttt{if},
\texttt{else}, \texttt{while}, \texttt{print}, the operators
\texttt{+ - * / = == != < <= > >=}, punctuation \texttt{( ) \{ \} ; ,}, and \texttt{\#}-style
line comments.

\begin{lstlisting}[language=Python]
KEYWORDS = {"if", "else", "while", "print"}

class Token:
    def __init__(self, kind, lexeme, line, col):
        self.kind, self.lexeme = kind, lexeme
        self.line, self.col = line, col

    def __repr__(self):
        return f"{self.kind}({self.lexeme!r})"


class LexError(Exception):
    def __init__(self, message, line, col):
        super().__init__(f"line {line}, col {col}: {message}")


class Lexer:
    def __init__(self, source):
        self.src = source
        self.pos = 0
        self.line, self.col = 1, 1

    def peek(self, ahead=0):
        i = self.pos + ahead
        return self.src[i] if i < len(self.src) else "\0"

    def advance(self):
        ch = self.src[self.pos]
        self.pos += 1
        if ch == "\n":
            self.line += 1
            self.col = 1
        else:
            self.col += 1
        return ch

    def scan(self):
        tokens = []
        while True:
            self._skip_whitespace_and_comments()
            if self.peek() == "\0":
                break
            tokens.append(self._scan_token())
        tokens.append(Token("EOF", "", self.line, self.col))
        return tokens

    def _skip_whitespace_and_comments(self):
        while True:
            c = self.peek()
            if c in " \t\r\n":
                self.advance()
            elif c == "#":
                while self.peek() not in ("\n", "\0"):
                    self.advance()
            else:
                return

    def _scan_token(self):
        line, col = self.line, self.col
        c = self.advance()

        # --- maximal munch for numbers ---
        if c.isdigit():
            lexeme = c
            while self.peek().isdigit():
                lexeme += self.advance()
            return Token("NUMBER", lexeme, line, col)

        # --- maximal munch for identifiers/keywords ---
        if c.isalpha() or c == "_":
            lexeme = c
            while self.peek().isalnum() or self.peek() == "_":
                lexeme += self.advance()
            kind = lexeme.upper() if lexeme in KEYWORDS else "IDENT"
            return Token(kind, lexeme, line, col)

        # --- maximal munch for two-character operators: try the ---
        # --- two-character form BEFORE falling back to one char ---
        two_char = {"<=": "LE", ">=": "GE", "==": "EQ", "!=": "NE"}
        two = c + self.peek()
        if two in two_char:
            self.advance()
            return Token(two_char[two], two, line, col)

        one_char = {
            "+": "PLUS", "-": "MINUS", "*": "STAR", "/": "SLASH",
            "=": "ASSIGN", "<": "LT", ">": "GT",
            "(": "LPAREN", ")": "RPAREN",
            "{": "LBRACE", "}": "RBRACE",
            ";": "SEMI", ",": "COMMA",
        }
        if c in one_char:
            return Token(one_char[c], c, line, col)

        raise LexError(f"unexpected character {c!r}", line, col)
\end{lstlisting}

The whole design is the maximal-munch loop in \texttt{\_scan\_token}: for each category,
keep consuming characters as long as the category's pattern is still satisfiable, and only
stop once extending the lexeme by one more character would break the match. The two-char
operator check is the clearest instance --- it explicitly tries the longer match
(\texttt{c + peek()}) before falling back to the single character, which is maximal munch
written out as an \texttt{if} statement.

\subsection{generated vs.\ hand-written lexers}

You didn't have to write that loop by hand. Tools like \texttt{lex} and its modern
descendant \texttt{flex} take a list of regular expressions with associated actions and
generate the scanning code (typically a table-driven DFA in C) for you. This is the
industrial approach: fast, consistent, and mechanically correct once your regular
expressions are right, because the tool is doing the DFA construction that you'd otherwise
have to do by hand.

The trade-off is control. A generated lexer is a black box that implements ``longest
match, first rule wins on ties'' uniformly; if your language needs something irregular ---
Python's indentation-sensitive blocks, for instance, which need a stack of indent levels
that isn't expressible as a single DFA --- you end up hand-patching the generated code
anyway. Robert Nystrom's \emph{Crafting Interpreters} makes the case for writing the lexer
by hand even for a real language (Lox): the code is short, it's exactly as fast as you need
for a tree-walking interpreter, and you get precise control over error messages and edge
cases like indentation or string interpolation. For a small DSL --- the kind you build once
and iterate on constantly --- hand-written usually wins for the same reason: you'll want to
tweak error messages and add odd little rules more often than you'll want raw scanning
throughput.

\begin{example}[Lexing a program by hand]
Take the following short program in the toy language above:
\begin{lstlisting}
x = 3
if x <= 10 {
  print x  # done
}
\end{lstlisting}
Running it through the lexer produces:

\begin{center}
\begin{tabular}{lll}
\hline
Token & Lexeme & Line \\
\hline
\texttt{IDENT}  & \texttt{x}     & 1 \\
\texttt{ASSIGN} & \texttt{=}     & 1 \\
\texttt{NUMBER} & \texttt{3}     & 1 \\
\texttt{IF}     & \texttt{if}    & 2 \\
\texttt{IDENT}  & \texttt{x}     & 2 \\
\texttt{LE}     & \texttt{<=}    & 2 \\
\texttt{NUMBER} & \texttt{10}    & 2 \\
\texttt{LBRACE} & \texttt{\{}    & 2 \\
\texttt{PRINT}  & \texttt{print} & 3 \\
\texttt{IDENT}  & \texttt{x}     & 3 \\
\texttt{RBRACE} & \texttt{\}}    & 4 \\
\texttt{EOF}    &                & 5 \\
\hline
\end{tabular}
\end{center}

Notice three things at once: \texttt{<=} survived as a single \texttt{LE} token (maximal
munch), \texttt{if} and \texttt{print} became keyword tokens rather than
\texttt{IDENT}s (the keyword table lookup), and the comment \texttt{\# done} disappeared
entirely, along with all the indentation --- it never influenced the token stream because
this toy language doesn't give whitespace any meaning.
\end{example}

\subsection{Exercises}
\begin{itemize}
  \item Lex the string \texttt{y != 4 \# stop} by hand, listing every token with its kind
    and lexeme, the way the worked example above does.
  \item Add support for \texttt{:=} as an alternative assignment operator, sitting
    alongside \texttt{=}. What does maximal munch have to say about the input
    \texttt{x := 3} versus \texttt{x : = 3}? Would your lexer accept the second one, and
    should it?
  \item Extend the \texttt{NUMBER} rule to accept a single decimal point (\texttt{3.14}).
    What ambiguity shows up if the input is \texttt{3.} followed immediately by
    \texttt{foo} with no space -- i.e., \texttt{3.foo}? How would you resolve it?
  \item Explain, in your own words, why the keyword check has to happen \emph{after}
    matching the full identifier pattern rather than before. Construct a concrete input
    that a ``keywords-first'' lexer would tokenize incorrectly.
  \item Modify \texttt{\_skip\_whitespace\_and\_comments} to support block comments
    delimited by \texttt{/* ... */} in addition to \texttt{\#} line comments. What happens
    if a block comment is never closed before end of file --- where should that error be
    reported?
  \item A generated lexer (flex-style) and the hand-written one above will tokenize any
    valid input identically. Describe one concrete situation where they would
    \emph{disagree} on how to report an error, and explain why.
\end{itemize}
